\subsection{Results}

\subsubsection{Encrypting a word}
The following sifted key was obtained from the experiment performed in
\cref{sec:bb84} in the case of no eavesdropping from
third parties:
\begin{equation*}
  1111 0001 1111 0000 0001 0110 01011 0 \, .
\end{equation*}
It is important to say that in real world, the sifted key needs to be further
processed to reduce third parties' eventual residual knowledge. This
is the famous privacy
amplification step.

As can be seen, the number of usable key bits (TODO:plural?) is 29.
The maximum size of the encryptable input for the OTP scheme is thus
composed by \SI{29}{\bit}.
TODO: put reference to where otp is defined.

The encryption was realized digitally, using
mainly a Python Jupyter notebook. Firstly a cell with both the key
and the data to encrypt,
chosen to be \lstinline[language=Python]|"INR"|, as an INRiM
truncation, was executed as in
\cref{lst:key_data}.
\begin{lstlisting}[language=python, caption=Cell with the initial data. , label=lst:key_data, float=h ]
token=0b111100011111000000010110010110
data="INR"
\end{lstlisting}
Notice that for the text, common ASCII encoding was used, resulting
in a byte of information per each character.
That is the reason why only three chars (so \SI{24}{\bit}) could be
encrypted with the given  \SI{29}{\bit} key.
Moreover, to make python understand well that the variable
\lstinline[language=Python]|token| is
a binary value, and so that the underlying information should be
understood as raw bits (s?),
the prefix \lstinline[language=Python]|0b| was used.

Then, \lstinline[language=Python]|data| was encoded using the
\lstinline[language=Python]|encode()| method,
specifying \lstinline[language=Python]|"utf-8"|,
a variable length encoding that has as a subset the ASCII encoding. This means
that python associates to the new encoded variable a now raw bit sequence. In
particular, the data associated to the data string is associated to
bytes according to \cref{tab:ascii}.
\begin{table}[h]
  \centering
  \begin{tabular}{cc}
    \toprule
    Character & Bitstring \\
    \midrule
    I & 0100 1001 \\
    N & 0100 1110 \\
    R & 0101 0010 \\
    \bottomrule
  \end{tabular}
  \caption{ASCII values for relevant characters.}
  \label{tab:ascii}
\end{table}
Python doesn't provide a way to perform logical computations like XOR
directly on encoded objects,
but requires another step firstly.
That consists in converting that binary stream into an integer number
and then display it
in the binary system. To do that,
\lstinline[language=Python]|int.from_bytes()| function
was used, followed by the \lstinline[language=Python]|format()|, with the
\lstinline[language=Python]|"b"| option, to indicate how should
python process the data. All this is shown compactly in \cref{lst:key_data}
\begin{lstlisting}[language=python, caption=Cell with the initial data. , label=lst:key_data , float=h ]
byte_encoded_data = data.encode("utf-8")
binary_encoded_data = int.from_bytes(byte_encoded_data)
binary_stream = format(binary_encoded_data, 'b')
\end{lstlisting}
\lstinline[language=Python]|binary_stream| is a variable that can be
XORed with the key, using
the \lstinline[language=Python]|^| operator. Before doing this, the
function \lstinline[language=Python]|cut_key()|, detailed in
\cref{lst:cut_key}, ensures that the key
has a bit length bigger than
the message one, raising an
\lstinline[language=Python]|AssertionError| if that is not the case.

\begin{lstlisting}[language=python, caption=The function that cuts the key making it compatible with the message. , label=lst:cut_key ,float=h ]
def cut_key(message, key):
    """
    Raises an error if the key is shorter than the message.
    """
    assert len(key) >= len(message), "Keylen < messlen"
    return key[:len(message)]
\end{lstlisting}

The final encryption step are shown in \cref{lst:final_enc}.

\begin{lstlisting}[language=python, caption=Final step. , label=lst:final_enc, float=h]
encrypted_mess = binary_stream ^ cut_key(binary_stream , token)
\end{lstlisting}

TODO: say how do Bob decrypt the thing.

\subsubsection{Encrypted communication}

A \texttt{keys} folder containing 2226 \texttt{.key} files, that are
files containing the
binary keys generated by the INRiM quantum key generator, was provided by
the researchers. The files were named with univoque hashed strings (TODO:true??)
and contained each \SI{32}{\byte} of data. For instance, running
\texttt{hexdump -C} on the file
\texttt{0a7431cf-6f5c-42ab-a3c5-d58b462cf963.key} produces the key
shown in \cref{lst:hexdump}.

\begin{lstlisting}[language=bash, caption=Output of \texttt{hexdump} command., label=lst:hexdump]
0000000 d7ce acc0 d082 1a22 8d1f b42c 1cb2 4509
0000010 ddbe 6478 9e5a 7dd4 16c7 b42c 5dc1 9dd6
0000020
\end{lstlisting}

TODO: finish